% /Users/pikachu/books/payments-book/src/parts/part1/ch06-regulation.tex
\chapter{Регуляторика: что требует закон}
\label{ch:regulation}

\begin{kaobox}[title=Что вы узнаете из этой главы]
  Предыдущие главы описывали, как устроен карточный платёж
  технически — протоколы, форматы, фазы транзакции. Но ни один
  из этих механизмов не существует в вакууме: каждый участник
  платёжной экосистемы работает в правовом поле, определяемом
  законами и нормативными актами. Здесь мы разберём основные
  регуляторные рамки — \law{161-ФЗ} и ключевые положения
  Банка России, — разберёмся, кому и какая лицензия нужна
  для работы с платежами, и кратко посмотрим на регуляторные
  подходы ближайших соседей: Казахстана, Беларуси и Армении.
\end{kaobox}

Полезно смотреть на регуляторику не как на приложение к технике,
а как на набор вопросов, которые закон задаёт продукту. Кто
именно переводит деньги? Кто имеет право маршрутизировать
сообщения и завершать расчёты? Нужна ли лицензия, где проходит
граница ответственности и какие операционные обязанности из этого
следуют?

\section{\law{161-ФЗ}: архитектура национальной платёжной системы}
\label{sec:regulation-161fz}

Федеральный закон от 27 июня 2011 года № 161-ФЗ
\enquote{О национальной платёжной системе} — базовый
нормативный акт, определяющий, кто и на каких условиях может
переводить деньги в России \sidecite{fz-161}. Для инженера,
работающего с платежами, \law{161-ФЗ} важен не столько
конкретными требованиями (они разбросаны по подзаконным
актам — положениям ЦБ), сколько терминологией и архитектурой
ролей: закон определяет, кто есть кто в платёжной системе,
и без понимания этих определений невозможно разобраться ни в одном
нормативном документе.

\subsection*{Оператор платёжной системы}

Центральное понятие \law{161-ФЗ} — \emph{платёжная система}:
совокупность организаций, правил и процедур, обеспечивающих
перевод денежных средств. Платёжную систему возглавляет
\emph{оператор платёжной системы} — организация, которая
определяет правила функционирования системы, контролирует
их соблюдение и несёт ответственность за её работоспособность.
В терминологии предыдущих глав оператор платёжной системы —
это карточная сеть: \scheme{Visa}, \scheme{Mastercard}, НСПК
(оператор платёжной системы Мир и СБП).

Оператором платёжной системы может быть кредитная организация
(банк), Банк России или организация, не являющаяся кредитной,
но зарегистрированная Банком России в реестре операторов
платёжных систем. Регистрация — не формальность: оператор
обязан утвердить правила платёжной системы, обеспечить управление
рисками, выбрать модель расчётов и выполнять требования
к бесперебойности \sidecite{fz-161}.

\subsection*{Три инфраструктурных роли}

\law{161-ФЗ} выделяет три вида организаций, обеспечивающих
инфраструктуру платёжной системы; каждая из них выполняет
свою функцию, и одна организация может совмещать несколько
ролей:

\begin{enumerate}
  \item \textbf{Операционный центр} (ОПКЦ — оператор услуг
        платёжной инфраструктуры в части операционных услуг) —
        обеспечивает обмен электронными сообщениями между
        участниками: передачу авторизационных запросов, ответов,
        reversal-ов. В терминах предыдущих глав — это процессинг,
        маршрутизирующий \iso{8583}-сообщения между эквайерами
        и эмитентами.
  \item \textbf{Платёжный клиринговый центр} (ОУИО — оператор
        услуг платёжной инфраструктуры в части услуг информационного
        обмена) — выполняет клиринг: приём, обработку
        и сопоставление клиринговых записей, расчёт нетто-позиций
        участников, формирование итоговых расчётных документов.
  \item \textbf{Расчётный центр} — обеспечивает завершение
        расчётов: фактическое перемещение денег между
        корреспондентскими счетами банков по итогам клиринга.
        Расчётным центром может быть Банк России или кредитная
        организация, назначенная оператором платёжной системы.
\end{enumerate}

НСПК совмещает все три роли для внутрироссийских транзакций:
выступает операционным центром (маршрутизирует авторизации),
платёжным клиринговым центром (обрабатывает клиринговые файлы)
и расчётным центром (производит нетто-расчёт между банками).
Именно поэтому с 2015 года все внутрироссийские карточные
транзакции — в том числе по картам \scheme{Visa}
и \scheme{Mastercard} — проходят через НСПК, а не через
зарубежные процессинговые центры карточных
сетей \sidecite{fz-161}.

\subsection*{Участники и операторы по переводу денежных средств}

Банки, фактически выполняющие переводы (эмитенты и эквайеры
в терминологии карточной индустрии), в \law{161-ФЗ} называются
\emph{операторами по переводу денежных средств}. Оператором
по переводу может быть только кредитная организация с лицензией
Банка России, Банк России или ВЭБ.РФ. Это ограничение означает,
что компания, не имеющая банковской лицензии, не может
самостоятельно переводить деньги — она может лишь оказывать
технические услуги (процессинг, платёжный шлюз), но сам перевод
должен выполнять лицензированный
банк \todo[inline]{Проверить: перечень организаций, имеющих статус оператора
  по переводу ДС, — ст.\,11 \law{161-ФЗ}}.

\begin{kaobox}[title=На практике]
  Если вы строите финтех-продукт в России и хотите переводить
  деньги, у вас два пути: получить банковскую лицензию (или лицензию
  НКО) и стать оператором по переводу самостоятельно — или заключить
  договор с банком и выступать в роли технического посредника
  (платёжного агрегатора, PSP), не касаясь денег клиентов напрямую.
  Первый путь даёт полный контроль, но требует капитала, комплаенса
  и отчётности перед ЦБ; второй — быстрее и дешевле, но ставит
  вас в зависимость от банка-партнёра.
\end{kaobox}

Итак, \law{161-ФЗ} задаёт словарь и архитектуру ролей: кто считается
оператором, какие инфраструктурные функции существуют и где проходит
правовая граница платёжной системы. Дальше этот каркас наполняется
конкретными ежедневными обязанностями уже на уровне положений
Банка России.

\section{Положения Банка России}
\label{sec:regulation-cbr-provisions}

\law{161-ФЗ} задаёт рамку, а конкретные требования прописаны
в подзаконных актах — положениях Банка России, каждое из которых
регулирует свой аспект платёжной деятельности. Для инженера,
работающего с карточными платежами, ключевыми являются три
положения: \law{382-П}, \law{683-П} и \law{266-П}.

\subsection*{\law{382-П}: защита информации при переводах}

Положение Банка России № 382-П устанавливает требования
к обеспечению защиты информации при осуществлении переводов
денежных средств \sidecite{cbr-382p}. Это основной \enquote{security}
документ для платёжных систем в России — российский ответ
на PCI DSS, хотя и существенно отличающийся по структуре
и подходу.

\law{382-П} обязывает операторов по переводу (банки),
операторов платёжных систем и операторов услуг платёжной
инфраструктуры выполнять комплекс мер по защите информации:
криптографическую защиту каналов связи, контроль доступа
к платёжной инфраструктуре, регистрацию инцидентов,
информирование Банка России о нарушениях безопасности
и регулярную оценку соответствия. Положение не заменяет
PCI DSS (его выполнение — требование карточных сетей,
а не ЦБ), но дополняет его российской спецификой: требованиями
к применению сертифицированных средств криптографической
защиты (СКЗИ), к отечественным стандартам шифрования (ГОСТ)
и к отчётности перед
регулятором \todo[inline]{Проверить: актуальная редакция 382-П —
  с учётом изменений, вступивших в силу после 2022 года}.

\subsection*{\law{683-П}: информационная безопасность в банках}

Положение № 683-П устанавливает обязательные для кредитных
организаций требования к обеспечению защиты информации
при осуществлении банковской
деятельности \sidecite{cbr-683p}. В отличие от \law{382-П},
которое фокусируется на переводах денежных средств,
\law{683-П} покрывает банковскую деятельность в целом —
включая дистанционное банковское обслуживание (ДБО),
мобильные приложения, интернет-банк и взаимодействие
с клиентами.

Для разработчика, работающего с банковскими системами,
\law{683-П} означает конкретные технические требования:
обязательное использование сертифицированного программного
обеспечения для обработки платежей, проведение анализа
уязвимостей, тестирование на проникновение, применение
ГОСТ-криптографии для защиты клиентских данных
и обязательное информирование ФинЦЕРТ (подразделение
Банка России, отвечающее за мониторинг киберинцидентов
в финансовой сфере) об инцидентах
безопасности \todo[inline]{Проверить: требования 683-П к ОУД4
  и сертификации ПО — проверить актуальные сроки вступления
  в силу}.

\subsection*{\law{266-П}: эмиссия платёжных карт}

Положение № 266-П регулирует порядок осуществления операций
с использованием платёжных карт — от выпуска карты до обработки
транзакций \todo[inline]{Источник: Положение ЦБ РФ № 266-П — добавить
  в bib}. Это положение определяет, какие операции допустимы
с платёжными картами, устанавливает требования к договорам
между участниками (эмитент и держатель, эквайер и мерчант),
регламентирует обязательное информирование клиента о совершённых
операциях и устанавливает порядок оспаривания.

Для инженера \law{266-П} важно в первую очередь тем, что оно
определяет обязательный набор реквизитов, которые должны
содержаться в документе по операции с платёжной картой
(электронном чеке): дату, сумму, валюту, реквизиты мерчанта,
код авторизации — по сути, то самое, что мы разбирали
в контексте data elements \iso{8583}.

\section{Обязательная эмиссия и эквайринг Мир}
\label{sec:regulation-mir-mandatory}

Отдельная и практически значимая тема — обязательность
карт Мир в российской платёжной инфраструктуре. Начиная
с 2017 года Банк России поэтапно вводил требования по приёму
и выпуску карт национальной платёжной системы Мир, и к настоящему
моменту они затрагивают практически всех участников
рынка \sidecite{fz-161}.

\subsection*{Обязательная эмиссия}

Все кредитные организации, осуществляющие эмиссию платёжных карт,
обязаны предлагать клиентам карты Мир — в том числе
для зачисления выплат из бюджета (зарплаты бюджетников, пенсии,
стипендии, социальные пособия). Получатели бюджетных средств
обязаны иметь карту Мир для зачисления этих
выплат \todo[inline]{Проверить: сроки и этапы обязательного перехода
  на Мир для бюджетных выплат — ст.\,30.5 \law{161-ФЗ},
  указания ЦБ}.

\subsection*{Обязательный эквайринг}

Торгово-сервисные предприятия с годовой выручкой свыше
порогового значения обязаны принимать к оплате карты Мир.
Порог периодически пересматривается: логика регулятора —
постепенное расширение охвата, начиная с крупнейших ретейлеров
и заканчивая малым
бизнесом \todo[inline]{Проверить: актуальный порог обязательного
  эквайринга Мир — проверить по \law{161-ФЗ} ст.\,30.5
  и постановлениям Правительства}.

\begin{kaobox}[title=Важно, colback=red!5, colbacktitle=red!20]
  Для разработчика, интегрирующего приём платежей, обязательный
  эквайринг Мир означает конкретное техническое требование:
  платёжная система мерчанта должна корректно обрабатывать карты
  Мир — распознавать IIN-диапазоны (2200--2204), маршрутизировать
  транзакции через НСПК и поддерживать специфику протокола Мир
  (собственные расширения \iso{8583}, идентификаторы продуктов,
  правила обработки contactless). Если ваш PSP или платёжный шлюз
  не поддерживает Мир — мерчант нарушает закон.
\end{kaobox}

\section{Лицензирование: кому нужна лицензия}
\label{sec:regulation-licensing}

Российское законодательство устанавливает разные требования
к лицензированию в зависимости от роли организации в платёжной
системе. Разобраться в них — необходимый шаг для любого,
кто планирует строить платёжный бизнес, а не только интегрировать
чужое API.

\subsection*{Кредитная организация (банк)}

Для осуществления банковских операций — в том числе
для переводов денежных средств, эмиссии карт и эквайринга —
необходима банковская лицензия, выдаваемая Банком России.
Банковская лицензия — самый \enquote{тяжёлый} вариант:
высокие требования к капиталу, масштабная отчётность,
обязательные нормативы достаточности, ликвидности
и резервирования \sidecite{fz-86}.

\subsection*{Небанковская кредитная организация (НКО)}

НКО — облегчённая форма кредитной организации, которая может
осуществлять отдельные банковские операции, определённые
лицензией, без полного набора банковских услуг. Для платёжной
индустрии наиболее релевантна расчётная НКО — организация,
имеющая право осуществлять переводы денежных средств без открытия
банковских счетов и связанные с ними операции. Многие
процессинговые компании и электронные кошельки в России работают
именно на лицензии расчётной НКО: она позволяет переводить деньги,
но не требует полного банковского капитала и инфраструктуры.

\subsection*{Оператор платёжной системы}

Регистрация в качестве оператора платёжной системы — процедура,
предусмотренная \law{161-ФЗ} для организаций, которые создают
и управляют собственными платёжными системами. Оператор обязан
направить в Банк России заявление о регистрации, приложив правила
платёжной системы, и получить запись в реестре операторов
платёжных систем. Банк России вправе отказать в регистрации,
если правила не соответствуют требованиям закона или организация
не способна обеспечить бесперебойность
работы \todo[inline]{Проверить: порядок регистрации оператора ПС —
  ст.\,15 \law{161-ФЗ}}.

\subsection*{ОПКЦ и ОУИО}

Оператор услуг платёжной инфраструктуры (ОПКЦ — операционный
центр, ОУИО — информационно-клиринговый центр) не нуждается
в банковской лицензии, но обязан соответствовать требованиям
оператора платёжной системы, в которой он работает,
и выполнять требования Банка России к защите информации
(\law{382-П}). На практике это означает, что процессинговая
компания, обрабатывающая авторизации для банка, может
не быть банком сама, но обязана пройти оценку соответствия
по \law{382-П} и обеспечить защиту обрабатываемых данных
на уровне, сопоставимом с банковским.

\begin{kaobox}[title=На практике]
  Если вы создаёте процессинговую платформу или платёжный шлюз
  в России, ваш правовой статус зависит от того, касаетесь ли
  вы денег: если платформа только маршрутизирует сообщения
  (ОПКЦ), банковская лицензия не нужна, но нужна оценка
  соответствия по \law{382-П} и договор с оператором платёжной
  системы. Если платформа открывает электронные кошельки
  или хранит средства клиентов — нужна лицензия НКО
  или банковская. Граница проходит по одному критерию: есть ли
  у вас обязательства перед клиентом по хранению или переводу
  \emph{его} денег.
\end{kaobox}

\section{Регулирование платежей в СНГ}
\label{sec:regulation-cis}

Платёжные системы не существуют в границах одной юрисдикции:
российские банки работают с контрагентами в странах СНГ,
финтех-компании запускают продукты на нескольких рынках
одновременно, а карты Мир принимаются в ряде
соседних стран. Краткий обзор регуляторных подходов
ключевых рынков СНГ позволит увидеть как общие черты
(все унаследовали советскую банковскую систему и перестраивали
её по схожим принципам), так и существенные различия.

\subsection*{Казахстан}

Казахстан регулирует платежи через Закон \enquote{О платежах
  и платёжных системах}, а надзор осуществляют Национальный банк
и Агентство по регулированию и развитию финансового рынка (АРРФР).
Казахстанский подход заметно либеральнее российского в части
допуска небанковских организаций: лицензию на оказание
платёжных услуг могут получить не только банки, но и платёжные
организации — юридические лица, не являющиеся банками, но имеющие
лицензию на предоставление платёжных
услуг \todo[inline]{Источник: Закон РК \enquote{О платежах и платёжных системах}}.

Национальная платёжная система Казахстана — межбанковская система
переводов денег (МСПД), управляемая Нацбанком; карточный процессинг
внутри страны обеспечивается Казахстанским центром межбанковских
расчётов (КЦМР). Казахстан активно развивает мгновенные платежи
и QR-оплату, а также принимает карты Мир на основании двусторонних
соглашений с НСПК.

\subsection*{Беларусь}

Регулирование платежей в Беларуси определяется Банковским кодексом
и рядом постановлений Национального банка Республики Беларусь.
Белорусская модель ближе к российской: оператором платёжной
системы может быть только Национальный банк или банк (кредитная
организация), а платёжная инфраструктура концентрируется вокруг
межбанковской расчётной системы BISS (Belarus Interbank Settlement
System), работающей по принципу
RTGS \todo[inline]{Источник: Банковский кодекс РБ, постановления НБРБ
  по платёжным системам}.

Национальная карточная система Беларуси — \enquote{БЕЛКАРТ},
управляемая ОАО \enquote{Банковский процессинговый центр}.
Карты Мир принимаются в Беларуси через инфраструктуру БЕЛКАРТ
на основании соглашения с НСПК, что позволяет российским
держателям карт использовать их в белорусских терминалах
и банкоматах.

\subsection*{Армения}

Центральный банк Армении регулирует платёжные системы
на основании Закона \enquote{О платёжных системах
  и платёжных услугах}. Армянский подход примечателен
ориентацией на европейские стандарты: закон во многом
вдохновлён директивой PSD2 Европейского союза и вводит
категории платёжных организаций, аналогичные европейским
(платёжные учреждения, учреждения электронных денег),
что отличает его от российской модели, основанной
на банковской
лицензии \todo[inline]{Источник: Закон РА \enquote{О платёжных системах
    и платёжных услугах}, влияние PSD2}.

Национальная платёжная система — ArCa, через которую проходят
внутренние карточные транзакции. Армения также принимает карты Мир
и активно развивает цифровые платежи, при этом сохраняя
относительно открытый режим для международных карточных сетей.

\begin{kaobox}[title=Важно, colback=red!5, colbacktitle=red!20]
  При запуске платёжного продукта в нескольких странах СНГ
  нельзя исходить из предположения, что регуляторные требования
  одинаковы. Различия начинаются с базового вопроса — кому
  нужна лицензия для перевода денег, — и продолжаются
  в требованиях к защите данных, к идентификации клиентов,
  к локализации инфраструктуры и к отчётности перед регулятором.
  Каждая юрисдикция требует отдельного правового анализа;
  \enquote{работает в России, значит, будет работать
    в Казахстане} — опасное допущение.
\end{kaobox}

\section*{Итоги главы}

\begin{itemize}
  \item \law{161-ФЗ} определяет архитектуру российской платёжной
        системы: оператор платёжной системы (задаёт правила),
        оператор по переводу денежных средств (банк, выполняющий
        переводы), оператор услуг платёжной инфраструктуры
        (ОПКЦ — операционный центр, ОУИО — клиринговый центр,
        расчётный центр).
  \item Три ключевых положения ЦБ: \law{382-П} (защита информации
        при переводах — аналог PCI DSS с российской спецификой),
        \law{683-П} (информационная безопасность в банках,
        сертификация ПО, ГОСТ-криптография), \law{266-П}
        (порядок операций с платёжными картами, реквизиты чека).
  \item Карты Мир обязательны к эмиссии для бюджетных выплат
        и к приёму для торговых предприятий сверх порога выручки;
        техническая поддержка Мир — не опциональное дополнение,
        а законодательное требование.
  \item Лицензирование зависит от роли: банковская лицензия —
        для переводов и эмиссии, лицензия НКО — облегчённый
        вариант, регистрация ОПКЦ/ОУИО — для процессинга
        без банковской лицензии, но с обязательной оценкой
        соответствия по \law{382-П}.
  \item Регуляторные подходы стран СНГ различаются: Казахстан
        допускает небанковские платёжные организации, Беларусь
        ближе к российской модели, Армения ориентируется
        на европейскую PSD2; при выходе на несколько рынков
        каждая юрисдикция требует отдельного анализа.
\end{itemize}
